%% CardioWatch — PSB 2027 SUBMISSION file (World Scientific ws-procs format).
%%
%% >>> SETUP: download the official PSB / World Scientific proceedings class from
%%     the PSB author instructions and place it beside this file (commonly
%%     ws-procs9x6.cls, with ws-procs.bst). If PSB's class name differs, change
%%     ONLY the \documentclass line below — the title/author/address/abstract/
%%     keywords/\bodymatter and all section/table/figure/\cite commands follow
%%     World Scientific conventions and transfer unchanged.
%% >>> Abstract and body are pulled verbatim from main.tex (single source of
%%     truth). The 6 \XX cells are filled from crossdevice_stats.json pre-submission.
%% >>> Keep <= 12 pages excluding title page, author list, and references.
\documentclass[]{ws-procs9x6}
\usepackage{graphicx}
\usepackage{booktabs}
\usepackage{amsmath}
\usepackage{siunitx}
\usepackage{xcolor}

\newcommand{\XX}{\textcolor{red}{\texttt{[run]}}}

\begin{document}

\title{Device-Agnostic Screening for Atrial Fibrillation on Consumer Wearables:
A Pre-Registered, Leakage-Safe Cross-Device Evaluation for Stroke Prevention}

\author{Urmi Shah}
\address{Department of Computer Engineering, San Jose State University,\\
San Jose, CA, USA\\ E-mail: urmimanishkumar.shah@sjsu.edu}

\begin{abstract}
Atrial fibrillation (AFib) is the most common sustained arrhythmia and a major,
largely preventable cause of ischemic stroke; a large share of cases are silent
and undiagnosed. Consumer smartwatches that record a single-lead
electrocardiogram (ECG) could surface these cases at population scale, but a
screening model is only useful if it works on the heterogeneous fleet of devices
people actually own---and a model trained on one recording device often collapses
on another. We treat this cross-device gap as the central obstacle to wearable
AFib screening and evaluate it with a pre-registered, leakage-safe protocol. On a
pre-registered, patient-grouped hold-out from a hospital 12-lead corpus
(CPSC~2018, $n{=}1{,}376$ records, Lead~I only), a lightweight device-agnostic
model that uses only heartbeat-timing features (RR~+~random forest) is
statistically indistinguishable from a deep convolutional--recurrent network
(CNN-LSTM) trained on the same data (AUC $0.935$ vs.\ $0.949$; DeLong $p{=}0.33$,
McNemar $p{=}0.40$, Holm-corrected), and this null is stable across $20$ alternate
hold-out seeds. Crucially, when both models are moved \emph{zero-shot} onto unseen
hardware---ambulatory Holter (MIT-BIH, Long-Term AF), single-lead AliveCor
(CinC~2017), and Apple~Watch---the timing model transfers (MIT-BIH AUC $0.907$)
while the deep model degrades; we quantify this head-to-head with paired DeLong
tests and patient-clustered confidence intervals. We frame the study as a
screening-methodology contribution: a reusable pre-registered evaluation of
cross-device robustness, not a new detector. An exploratory onset-prediction
analysis (104 patients) yields only modest, sub-clinical discrimination
(AUC~$0.66$) and is reported as such. All headline numbers come from a single
pre-registered run and are released for reproduction.
\end{abstract}

\keywords{atrial fibrillation; single-lead ECG; wearable screening; cross-device
generalization; device-agnostic features; pre-registered evaluation; stroke prevention}

\bodymatter
\section{Introduction}
Atrial fibrillation is the most common sustained arrhythmia and a leading cause
of ischemic stroke, yet a large fraction of cases are paroxysmal, silent, and
undiagnosed until a stroke occurs. The clinical benefit of finding AFib is
concrete and time-sensitive: timely anticoagulation substantially reduces stroke
risk~\cite{perez2019appleheart}. Because so many cases are asymptomatic, the
bottleneck is not treatment but \emph{detection at scale}, and consumer
wearables---smartwatches that record a single-lead ECG (for example, Apple~Watch
Series~4 and later)---put an opportunistic screening instrument in tens of
millions of hands. The AFib signature (absent P~waves and irregular beat-to-beat
timing) is visible in a single lead, so single-lead screening is feasible in
principle.

The obstacle we address is a \emph{deployment} problem specific to screening.
A screening model does not run on one curated device; it must run, unchanged, on
whatever hardware a population already wears---different sampling rates, gains,
electrode geometries, and noise profiles. Models trained on one recording device
frequently degrade sharply on another because they latch onto hardware-specific
waveform characteristics rather than the underlying cardiac pattern; we call this
the \emph{cross-device gap}. For population screening, robustness to this gap is
not a secondary metric---it determines who actually gets caught.

This paper asks: for single-lead AFib screening that must generalize across the
consumer-device fleet, how large is the cross-device gap, and does a simple
device-agnostic representation close it? Our contributions, in the order we
consider them most important for a clinical-screening venue, are:
\begin{enumerate}
  \item A \textbf{pre-registered, leakage-safe evaluation protocol} for wearable
  AFib screening---frozen patient-grouped hold-out, frozen thresholds, and a
  pre-declared family of paired significance tests---directly responsive to the
  documented leakage/reproducibility problems in medical machine
  learning~\cite{kapoor2023leakage}.
  \item A \textbf{zero-shot cross-device evaluation} that scores a device-agnostic
  timing model and a deep waveform model on the \emph{same} unseen recordings from
  four hardware types (hospital 12-lead, ambulatory Holter, single-lead AliveCor,
  and Apple~Watch), with paired DeLong tests and patient-clustered confidence
  intervals.
  \item A \textbf{controlled same-data comparison} showing that, on matched
  hospital data, the timing model is not measurably worse than the deep model, a
  null result that is stable across seeds---bounding, rather than overselling, the
  ``simple versus deep'' question.
  \item An \textbf{honest exploratory study} of AFib \emph{onset} prediction from
  pre-onset sinus rhythm, reported with its limitations.
\end{enumerate}
We scope the work to AFib detection and stroke prevention; it does not address
heart-attack (myocardial infarction) prediction, a different condition with a
different ECG signature.

\section{Related Work}
Single-lead AFib detection has been studied since the PhysioNet/CinC 2017
challenge on short AliveCor recordings~\cite{clifford2017af}, and large
prospective work established that smartwatch screening surfaces undiagnosed
AFib~\cite{perez2019appleheart}. Recent deep models report very high in-domain
accuracy, and---directly relevant to us---Ben-Moshe et
al.~\cite{benmoshe2024rawecgnet} show that a carefully trained raw-ECG network can
be made to \emph{generalize} across leads and cohorts, and argue that raw-waveform
models can outperform hand-crafted RR-interval features. We do not dispute this:
with cross-device training and enough data, learned representations can close the
gap. Our claim is narrower and complementary. First, the realistic screening
constraint is \emph{zero-shot} transfer---a deployed model rarely gets retrained
per device---and under that constraint we find timing features transfer without
any cross-device training. Second, we make the comparison on \emph{matched} data
with a pre-registered paired test, rather than comparing models trained on
different corpora. Bahrami~Rad et al.~\cite{bahramirad2024crowdsourced} similarly
report that timing-based features transfer across AliveCor and Apple~Watch better
than learned waveform features; we add a same-data statistical comparison and a
multi-cohort, clustered-CI cross-device evaluation.

Methodologically, our design responds to the leakage-driven reproducibility
crisis in ML-based science~\cite{kapoor2023leakage}: we pre-register the hold-out
and tests before training. Onset prediction has a long history beginning with the
PAF Prediction Challenge~\cite{moody2001paf}; we revisit it on continuous Holter
data with patient-grouped evaluation and treat it as exploratory.

\section{Data and Methods}

\subsection{Datasets}
We use five public or personally exported sources. CPSC~2018 (within the
PhysioNet/CinC~2020 corpus~\cite{alday2020classification}) provides hospital
12-lead ECGs, from which we take Lead~I only. PhysioNet/CinC~2017 provides
single-lead AliveCor recordings~\cite{clifford2017af}. The MIT-BIH Atrial
Fibrillation Database (\texttt{afdb}) and the Long-Term AF Database
(\texttt{ltafdb}) provide annotated ambulatory Holter
recordings~\cite{moody1983afdb,petrutiu2007ltaf,goldberger2000physionet}. A small
set of personal Apple~Watch exports (54 recordings, four volunteers, one confirmed
AFib) is used only for real-world validation; its labels are the watch's own
classifier output, not independent clinical adjudication. A separate tabular
clinical risk dataset is used for an additional analysis (Section~\ref{sec:extra})
and is reported as a distinct cohort.

\subsection{Models}
The \textbf{device-agnostic model} (RR~+~RF) detects AFib from heartbeat-timing
features only---coefficient of variation, RMSSD, pNN50, sample entropy, and
related statistics of the RR-interval series---fed to a random forest. Because
these features measure timing rather than waveform shape, they are insensitive to
device sampling rate and gain. The \textbf{deep model} (CNN-LSTM) is a three-block
1D convolutional encoder followed by a two-layer LSTM and a linear head, trained
on 10-second Lead~I windows.

\subsection{Pre-registered evaluation protocol}
To avoid comparing models trained on different data, and to guard against the
test-set leakage that undermines reproducibility in medical
ML~\cite{kapoor2023leakage}, we fixed the analysis before training. We created a
deterministic, label-stratified hold-out of CPSC records (20\%, fixed seed) after
an explicit one-record-per-patient check, and excluded these records from
\emph{both} models' training. The device-agnostic model and a CPSC-only deep model
were then scored on the identical hold-out. Decision thresholds were frozen on
training/validation data, never on the hold-out. The primary discrimination test
is the paired DeLong test on AUC~\cite{delong1988comparing,sun2014fast}; the
primary threshold-level test is McNemar's test on correctness vectors. A single
pre-registered Holm family~\cite{holm1979simple} (DeLong plus fixed-threshold
McNemar across the model-versus-model pairs) controls multiple comparisons.
Class imbalance during training was handled with SMOTE applied \emph{inside} each
cross-validation fold~\cite{chawla2002smote}. A pre-registered amendment adds the
cross-device family below; the same-data primary is unchanged.

\subsection{Zero-shot cross-device evaluation}
The screening question is answered by moving the \emph{already-trained} models,
with no retraining or fine-tuning, onto unseen hardware. For each external cohort
we score RR~+~RF and the CPSC-only CNN-LSTM on the \emph{identical} 10-second
Lead~I windows, so the comparison is truly paired: long Holter records are cut
into consecutive 10-second windows and each window is labeled by the active rhythm
at its midpoint. Per cohort with both classes present we run a paired DeLong test
of RR~+~RF versus CNN-LSTM(CPSC); these cross-device tests form a single
pre-registered Holm family, separate from the same-data family. Every external AUC
is reported with a \textbf{patient-/record-clustered bootstrap} 95\% CI (resampling
whole patients, since windows within a patient are correlated). The
data-asymmetric combined CNN-LSTM is scored only on cohorts it never saw during
training (it is omitted for CinC~2017 to avoid train-on-test leakage).

\subsection{Seed robustness (secondary)}
The primary same-data null uses the pre-registered seed. As a labeled secondary
sensitivity analysis, we rerun the entire hold-out\,$\rightarrow$\,train\,$\rightarrow$\,paired-eval
chain over 20 alternate hold-out seeds and report the distribution of $\Delta$AUC
and DeLong $p$. This does not replace the pre-registered primary.

\subsection{Exploratory onset prediction}
Separately, we test whether imminent AFib can be predicted from sinus rhythm. From
continuous Holter records we label 30-second sinus windows as \emph{pre-onset}
(onset within 30~minutes) or \emph{control} (no onset within 60~minutes), excluding
windows already in AFib. Features combine RR/HRV statistics with
neurokit2~\cite{makowski2021neurokit} atrial-ectopy and expanded HRV markers, each
z-scored against the patient's own control-window baseline. The pre-registered
primary is the neurokit-feature variant; RR-only and P-wave variants are reported
as labeled sensitivity analyses. Evaluation uses patient-grouped cross-validation
pooled over \texttt{afdb} and \texttt{ltafdb}.

\section{Results}

\subsection{Cross-device generalization (screening centerpiece)}
Table~\ref{tab:crossdevice} reports the zero-shot head-to-head. Without any
retraining the device-agnostic model transfers across hardware, reaching AUC
$0.907$ on ambulatory Holter (MIT-BIH; 28{,}104 windows from 25 patients), while
the CPSC-only deep model degrades off its training device. The gap is quantified
per cohort by paired DeLong on identical windows and by patient-clustered CIs;
this---not raw in-domain accuracy---is the property that matters for screening on
a heterogeneous device fleet. The Apple~Watch row has a single confirmed positive,
so it is a plausibility check (accuracy dominated by specificity, sensitivity
effectively unmeasured; device-classifier labels), not a discrimination claim.

\begin{table}[t]
\centering
\caption{Zero-shot cross-device head-to-head on identical 10\,s Lead~I windows.
AUCs carry patient-/record-clustered bootstrap 95\% CIs; DeLong $p$ is
Holm-corrected across the cross-device family. Cells marked \XX{} are filled from
the pre-registered run (\texttt{crossdevice\_stats.json}).}
\label{tab:crossdevice}
\begin{tabular}{llcccc}
\toprule
Cohort & Device & RR~+~RF AUC & CNN-LSTM(CPSC) AUC & $\Delta$AUC & DeLong $p$(Holm) \\
\midrule
MIT-BIH \texttt{afdb} & Holter 250\,Hz     & 0.907 & \XX & \XX & \XX \\
Long-Term AF \texttt{ltafdb} & Holter 128\,Hz & \XX & \XX & \XX & \XX \\
CinC 2017 & AliveCor 300\,Hz               & \XX & \XX & \XX & \XX \\
Apple~Watch$^{\dagger}$ & Wearable 512\,Hz  & \multicolumn{4}{c}{plausibility only (1 positive; acc.\ $0.907$, 95\% CI $0.80$--$0.96$)} \\
\bottomrule
\end{tabular}
\end{table}

\subsection{Controlled same-data comparison}
On the pre-registered hold-out (Table~\ref{tab:ecg}), the device-agnostic timing
model and the CPSC-only deep model are statistically indistinguishable: the AUC
difference is $-0.013$ (95\% CI $-0.032$ to $0.006$), DeLong $p{=}0.33$ and McNemar
$p{=}0.40$ after Holm correction. Across 20 alternate hold-out seeds this null is
stable ($\Delta$AUC and DeLong $p$ distributions in
Table~\ref{tab:ecg}, footnote). The combined deep model, additionally trained on
wearable data, reaches AUC $0.975$ and is significantly higher, but this is a
deployment comparison with more and different training data, not a same-data test.

\begin{table}[t]
\centering
\caption{ECG models on the pre-registered CPSC hold-out ($n{=}1{,}376$, Lead~I).
Thresholds frozen on non-hold-out data. Seed-robustness summary
(\texttt{seed\_robustness.json}) is reported alongside as a secondary sensitivity
analysis.}
\label{tab:ecg}
\begin{tabular}{lccccc}
\toprule
Model & AUC & Recall & Specificity & F1 & Threshold \\
\midrule
RR~+~RF (device-agnostic) & 0.935 & 0.803 & 0.879 & 0.679 & 0.40 \\
CNN-LSTM (CPSC only) & 0.949 & 0.959 & 0.832 & 0.701 & 0.40 \\
CNN-LSTM (combined, deployment) & 0.975 & 0.918 & 0.941 & 0.837 & 0.40 \\
Random baseline & 0.502 & 0.508 & 0.511 & 0.269 & 0.50 \\
\bottomrule
\end{tabular}
\end{table}

\subsection{Exploratory onset prediction}
Predicting AFib before it begins is a harder, separate problem. Pooled across two
Holter cohorts (104 patients, patient-grouped cross-validation), personalized
timing and atrial-ectopy features give modest, consistent pre-onset discrimination
(neurokit variant, window AUC $0.66$). At about $0.5$ false alarms per hour the
model warns of a minority of onsets a median of $\sim$1--2 minutes ahead.
Per-patient normalization was essential. We report this as a real but sub-clinical
signal, not a deployable early-warning system; the RR-only ($0.606$) and P-wave
variants are sensitivity analyses.

\subsection{Additional analysis: tabular clinical risk (separate cohort)}
\label{sec:extra}
For completeness we also train tabular clinical risk models (random forest, AUC
$0.954$; gradient boosting, AUC $0.927$) on a distinct patient population. These
are \emph{not} combined with the ECG models: the cohorts do not overlap, so any
fused score would be unvalidated. We therefore report them as a separate cohort
and treat clinical$+$ECG fusion as out of scope until same-patient paired data are
available.

\section{Discussion}
The screening-relevant result is the cross-device one: a simple, interpretable
timing representation transfers zero-shot across recording hardware, while a deep
waveform model trained on a single device does not. For opportunistic AFib
screening on the consumer-device fleet---where the model cannot be retrained per
watch---this favors the device-agnostic representation as a default. This does not
contradict work showing raw-ECG networks \emph{can} generalize with cross-device
training~\cite{benmoshe2024rawecgnet}; it says that under the zero-shot constraint,
and on matched data, the simple representation is competitive and more portable.
The combined deep model scores higher only with extra, heterogeneous training
data---a statement about data, not architecture.

Several limitations bound our claims. The Apple~Watch validation is small, has one
positive case, and relies on device-generated labels, so it supports plausibility
rather than sensitivity. Holter cross-device intervals are computed over windows
that are correlated within patients; we therefore report patient-clustered CIs
rather than naive per-window CIs. The clinical and ECG models use different
populations and are not a cross-modal ablation. The onset analysis uses cohorts of
selected AFib patients and shows only modest discrimination; richer atrial markers
(such as P-wave morphology), longer pre-onset context, and purpose-built data are
the natural next steps.

\section{Conclusion}
For single-lead AFib screening on consumer wearables, robustness to the
cross-device gap---not in-domain accuracy---is the property that determines who
gets caught. A device-agnostic timing model matches a deep model on matched data,
is stable across seeds, and transfers zero-shot across hardware, making it a
strong, interpretable default for stroke-prevention screening. We contribute a
pre-registered, leakage-safe protocol and a multi-cohort cross-device evaluation
that others can reuse. Predicting AFib before it starts remains hard: our
personalized features yield a real but sub-clinical early-warning signal that
motivates future work rather than a deployable alarm.

\section*{Reproducibility}
All reported numbers come from one pre-registered run. Code, the analysis plan
(including the cross-device amendment), the hold-out manifest checksum, and the
compact result files are available in the project repository; a versioned snapshot
will be archived with a DOI.

\bibliographystyle{ws-procs}
\bibliography{references}

\end{document}
